\section{Unsmoothed windows and an averaged common raw remainder}
\label{sec:unsmoothed-window-remainders}
The transport identity compares two raw corrections but does not bound
one of them. Positivity gives a different estimate: it bounds the
integral of the common correction over an actual observation window.
The window is not convolved with noise. Its side lengths must exceed
the resolution of a proved Fourier region. We keep this restriction
explicit, in particular in the three integer coordinates.

Put $X_{n,R}=U_{n,R}/\sqrt n$, $U_{n,R}=J_{n,R}-n\bar G_R$, and set
\begin{equation}\label{eq:window-Fourier-error}
 e=\frac3{280},\quad v_0=\frac9{175},\qquad
 \mathcal E_n(a)=M_an^{-e}\sqrt{\log(2+n)}+V_an^{-v_0}.
\end{equation}
The centered pushforward of the marked measure in
\eqref{eq:marked-measure} is denoted by $\lambda_{n,R}^{[k],a}$.
Its mass is $\alpha_a$. Theorem~\ref{thm:adaptive-cross-fixed-count}
gives the integrated characteristic error $C\mathcal E_n(a)$ on
$\mathcal U_n^\sharp$. Frequencies here are rescaled; a frequency $v$
corresponds to the physical frequency $v/\sqrt n$.

\subsection{Order-preserving interval approximation}
We give the approximation argument, including its values at endpoints.
These values matter because three coordinates of the record are discrete.
The construction is the classical Beurling--Selberg interval construction;
see \cite[Introduction]{CVExt}. Only the elementary facts proved below
are used, not an extremality assertion.

\begin{lemma}[Bandlimited envelopes, including endpoints]
\label{lem:window-interval-envelopes}
For an interval $I$ of length $l>0$ and an angular bandwidth $b>0$,
there are real continuous integrable functions $L_I,U_I$ with Fourier
support in $[-b,b]$ such that
\begin{equation}\label{eq:interval-envelopes}
 L_I\le\one_I\le U_I,\quad 0\le U_I\le3,\quad
 \int U_I=l+\frac{2\pi}{b},\quad
 \int L_I=l-\frac{2\pi}{b}.
\end{equation}
Either convention at either endpoint is allowed in the same inequalities.
The difference $E_I=U_I-L_I$ is nonnegative and has integral $4\pi/b$.
The Fourier convention in this statement is
$\widehat f(v)=\int e^{-ivx}f(x)\dd x$.
\end{lemma}
\begin{proof}
Use the continuous extensions of
\[
 K(x)=\left(\frac{\sin\pi x}{\pi x}\right)^2,\qquad
 H(x)=\sum_{m\in\Z}\operatorname{sgn}(m)K(x-m)+2xK(x),
 \quad \operatorname{sgn}(0)=0.
\]
The identity $\sum_{m\in\Z}K(x-m)=1$ follows from the partial-fraction
identity for $\pi^2/\sin^2\pi x$; the limiting integer cases agree.
For $x>0$ not an integer,
\[
 1-H(x)=K(x)+2\left(\frac{\sin\pi x}{\pi}\right)^2
                  \left(\sum_{j\ge1}\frac1{(x+j)^2}-\frac1x\right).
\]
The decreasing-integrand comparisons give
$1/(x+1)\le\sum_{j\ge1}(x+j)^{-2}\le1/x$.
Since $1/(x+1)\ge1/x-1/x^2$, it follows that
$|1-H(x)|\le K(x)$. Oddness proves
$|\operatorname{sgn}(x)-H(x)|\le K(x)$ everywhere.
Moreover $H-\operatorname{sgn}$ is an odd integrable function.

In the Fourier convention $e^{-2\pi i\xi x}$, $K$ has transform
$(1-|\xi|)_+$: write the sinc as the inverse transform of an interval
and take its square. Each translate of $K$, and $xK$ as a tempered
distribution, has support in $[-1,1]$ after Fourier transformation.
The partial sums defining $H$ converge as tempered distributions,
because their absolute values are bounded by $\sum_m K(x-m)=1$.
Thus $H$ has the same distributional Fourier support.

Write the interval endpoints as $c<d$, and put $\delta=b/(2\pi)$.
Set
\[
 U_I(x)=\tfrac12\{H(\delta(x-c))-H(\delta(x-d))
                       +K(\delta(x-c))+K(\delta(x-d))\},
\]
and change the two plus signs before $K$ to minus signs to obtain $L_I$.
The sign comparison proves the bounds off the endpoints. At $x=c$,
$U_I(c)=(H(\delta(d-c))+1+K(\delta(d-c)))/2\ge1$ and
$L_I(c)=(H(\delta(d-c))-1-K(\delta(d-c)))/2\le0$;
the other endpoint is the same. This proves both endpoint conventions.
The pointwise errors are bounded by the sum of the two $K$ terms,
so $0\le U_I\le3$. The odd integrable error of $H$ has integral
zero, and $\int K=1$. Hence the stated integrals follow. The
support statement follows by translation and dilation of the
compactly supported distributions above; the final functions are
integrable, so their Fourier transforms are continuous. In particular
Fourier inversion and integration against a finite measure are valid.
\end{proof}

For $B=\prod_{i=1}^d I_i$, where $d=3$ or $4$, write
$|B|=\prod_i l_i$ and $\mathcal Q(\mathbf b)=\prod_i[-b_i,b_i]$.
If $b_il_i\ge1$, tensorization gives envelopes with support in
$\mathcal Q(\mathbf b)$ and
\begin{equation}\label{eq:box-envelope-norms}
 L_B\le\one_B\le U_B,\quad
 \|L_B\|_1+\|U_B\|_1\le C_d|B|,\quad
 \int(U_B-\one_B)+\int(\one_B-L_B)
               \le C_d|B|\sum_i(b_il_i)^{-1}.
\end{equation}
Here the minorant is not the product of the minorants, which can
have the wrong sign. Precisely, set
\begin{equation}\label{eq:box-lower-envelope}
 U_B=\prod_i U_{I_i},\qquad
 L_B=U_B-\sum_i E_{I_i}\prod_{j\ne i}U_{I_j}.
\end{equation}
To verify the lower bound, telescope
$\prod_i U_{I_i}-\prod_i\one_{I_i}$ and use
$0\le U_{I_i}-\one_{I_i}\le E_{I_i}$ and $\one_{I_j}\le U_{I_j}$.
Integration of these nonnegative error bounds proves
\eqref{eq:box-envelope-norms}. Products here involve different variables,
so they do not increase a coordinate bandwidth.

\subsection{An actual box denominator, without a smoothing event}
Let $\mathbf h=(h_1,\ldots,h_4)$, $0<h_i\le1$, and define
\begin{equation}\label{eq:physical-window-box}
 B(\xi,\mathbf h)=\prod_i[\xi_i-h_i,\xi_i+h_i],\quad
 V(\mathbf h)=\prod_i2h_i,\quad
 \mathcal B_{n,R}(\xi,\mathbf h)=n\bar G_R+\sqrt n B(\xi,\mathbf h).
\end{equation}
In the last expression the first three coordinates are intersected
with $\Z^3$. Write $\mathfrak m$ for counting measure times Lebesgue
measure on the mixed record space, and put
\begin{equation}\label{eq:window-mixed-Gaussian}
 \mathfrak G_{n,R}(\mathcal B)=n^{-2}\int_{\mathcal B}
       g_{D_R}\left(\frac{x-n\bar G_R}{\sqrt n}\right)\dd\mathfrak m(x).
\end{equation}
The Gaussian density has the standard probability normalization.

\begin{theorem}[Unsmoothed window probabilities at a prescribed count]
\label{thm:unsmoothed-window-denominator}
Fix $A<\infty$. Suppose $|\xi|\le A$, $b_ih_i\ge1$,
$0<b_i\le\sqrt n/2$, and
$\mathcal Q(\mathbf b)\subset\mathcal U_n^\sharp$.
Then, uniformly in $R,n,k,a$ and these boxes,
\begin{align}
 \left|\mu_{n,R}^{[k],a}(\mathcal B)
                      -\alpha_a\mathfrak G_{n,R}(\mathcal B)\right|
 \le C_A V(\mathbf h)
       \left(\mathcal E_n(a)+M_a\sum_i(b_ih_i)^{-1}\right).
 \label{eq:unsmoothed-window-denominator}
\end{align}
For $a\ge0$, the factor $M_a$ in the second term may be replaced by
$\alpha_a$. If $\alpha_a>0$, $a\ge0$, and
\begin{equation}\label{eq:window-relative-budget}
 \epsilon_n(a,\mathbf b,\mathbf h)
       =\mathcal E_n(a)/\alpha_a+\sum_i(b_ih_i)^{-1}\longrightarrow0,
\end{equation}
then the ratio of the probability to its displayed Gaussian mass
converges to one uniformly on these central windows.
\end{theorem}
\begin{proof}
First suppose $a\ge0$, so $\lambda_{n,R}^{[k],a}$ is a positive
measure. Integrate the two envelopes against it. Their transforms
are supported in $\mathcal Q(\mathbf b)$ and bounded by $C V(\mathbf h)$,
by \eqref{eq:box-envelope-norms}. Fourier inversion and the inherited
integrated characteristic estimate give
\[
 \left|\int F\dd\lambda_{n,R}^{[k],a}
               -\alpha_a\int Fg_{D_R}\dd x\right|
       \le C V(\mathbf h)\mathcal E_n(a),\qquad F=L_B,U_B.
\]
This integrates an entire test function against the actual discrete
coordinates; no density with respect to four-dimensional Lebesgue
measure is assigned to that law. Uniform ellipticity bounds
$\|g_{D_R}\|_\infty$. The envelope excess then proves the same
estimate with $\int_Bg_{D_R}$ as Gaussian mass, with an additional
$C\alpha_a V(\mathbf h)\sum_i(b_ih_i)^{-1}$.

We now compare that continuous Gaussian mass with the mixed mass
in \eqref{eq:window-mixed-Gaussian}. The grid spacing is $n^{-1/2}$
in precisely the first three standardized coordinates. In one
coordinate, assign to each included grid point its interval of that
length. The symmetric difference from the target interval has length
at most $2n^{-1/2}$. On full cells use the uniform first-derivative
bound of the Gaussian. Iterating in the three grid coordinates gives
\begin{equation}\label{eq:window-grid-error}
 \left|\mathfrak G_{n,R}(\mathcal B)-\int_Bg_{D_R}\dd x\right|
 \le C_A V(\mathbf h)\sum_{i=1}^3(\sqrt n h_i)^{-1}.
\end{equation}
The conditions imply $\sqrt n h_i\ge2$, so all intermediate grid
volumes are comparable to the interval lengths. Since $b_i\le\sqrt n/2$,
the grid error is absorbed by the stated budget. This explicitly pays
for endpoints and the lattice convention.

For complex $a$, decompose its real and imaginary parts into positive
and negative parts. Their supremum and variation norms have sum at
most a fixed multiple of $M_a+V_a$, by contraction of BV under positive
part. Their masses have sum at most $C M_a$. Add the positive-measure
estimates to obtain \eqref{eq:unsmoothed-window-denominator}.
For the last assertion, uniform ellipticity and $|\xi|\le A$, $h_i\le1$
give $\mathfrak G_{n,R}(\mathcal B)\ge c_A V(\mathbf h)$ for all
sufficiently large $n$ under the vanishing budget. Divide the positive
estimate by that lower bound times $\alpha_a$.
\end{proof}

\begin{corollary}[A concrete rare-window law]
\label{cor:concrete-four-window}
Take $a=s_R$, so $\alpha_a=1$ and the original section probability
is recovered. Set $\rho=67/1400$ and $h_i=n^{-1/25}$.
Uniformly for $|\xi|\le A$,
\begin{equation}\label{eq:concrete-four-window}
 \nu_R^*\{X_{n,R}\in B(\xi,\mathbf h)\}
  =16n^{-4/25}g_{D_R}(\xi)
       \left(1+O_A(n^{-11/1400}\sqrt{\log(2+n)})\right).
\end{equation}
The event uses the original record, integer intervals in its first
three coordinates, and a time interval in its fourth. Each physical
half-width is $n^{23/50}$. No convolution is made in its definition.
\end{corollary}
\begin{proof}
Choose $b_i=n^\rho/4$. The product frequency box lies in the retained
isotropic ball of radius $2n^\rho$. Here
$\rho-1/25=11/1400>0$, $e=15/1400$, and $v_0>e$.
Theorem~\ref{thm:unsmoothed-window-denominator} and the Gaussian
first-derivative bound give the conclusion, since replacing the
Gaussian mass by its value at the center times volume costs relative
$O(n^{-1/25})$. The grid error is smaller. The Fourier supports are
rescaled $b_i$ and physical $b_i/\sqrt n=n^{-633/1400}/4$.
\end{proof}

\subsection{A size bound for the common signed correction}
\begin{theorem}[The common correction averaged over an actual window]
\label{thm:averaged-common-raw-remainder}
In addition to the hypotheses of
Theorem~\ref{thm:unsmoothed-window-denominator}, suppose the original
marked weight is finite-record admissible, and take the same finite
extraction and $L_n=\lceil\lambda n\rceil$, $\lambda>c_*^{-1}+1$,
as in Section~\ref{sec:coherent-raw-cutoff}. For any of the three
cutoffs $\chi_n^\sharp,\chi_n^{\rm box},\chi_n^\natural$, and any
fixed $P>0$,
\begin{align}
 \frac{n^2}{\mathfrak m(\mathcal B)}
 \left|\int_{\mathcal B}\mathcal R_\chi^{a,L_n}\dd\mathfrak m\right|
 \le C_A\left(\mathcal E_n(a)+M_a\sum_i(b_ih_i)^{-1}
                                      +M_an^{-P}\right).
 \label{eq:averaged-common-raw-remainder}
\end{align}
Thus for the concrete windows in \eqref{eq:concrete-four-window} the
common correction itself, rather than only a change between two
cutoffs, has a vanishing normalized signed window integral.
The absolute value is outside the integral.
\end{theorem}
\begin{proof}
Every cutoff is between zero and one, equals one on the original
central ball, and is supported in the retained Fourier union.
Mixed inversion and the integrated characteristic estimate imply
\begin{equation}\label{eq:window-smoothed-density}
 \sup_x\left|n^2(K_\chi*\mu_{n,R}^{[k],a})(x)
          -\alpha_ag_{D_R}((x-n\bar G_R)/\sqrt n)\right|
       \le C\mathcal E_n(a)+C_P M_an^{-P}.
\end{equation}
The omitted Gaussian frequencies are bounded by uniform ellipticity;
the first three physical supports are inside one torus chart.
On the central box, $N\le L_n$ for all sufficiently large $n$, so
$\mu^{a,L_n}(\mathcal B)=\mu_{n,R}^{[k],a}(\mathcal B)$ exactly.
The exact identity \eqref{eq:exact-raw-cutoff-transport} gives
\[
 \int_{\mathcal B}\mathcal R_\chi^{a,L_n}
  =\mu_{n,R}^{[k],a}(\mathcal B)
     -\int_{\mathcal B}K_\chi*\mu_{n,R}^{[k],a}
     +\int_{\mathcal B}K_\chi*\mathsf T^{a,L_n}.
\]
All integrals use $\mathfrak m$. Apply the preceding theorem to the
first term, \eqref{eq:window-smoothed-density} to the second, and
count-coordinate kernel separation to the last. The third-coordinate
support gap is of order $n$, and each of the three kernels gives
any fixed inverse power on that gap. Finally
$\mathfrak m(\mathcal B)\asymp n^2V(\mathbf h)$ by the same grid
count used in \eqref{eq:window-grid-error}. This proves the estimate.
No derivative norm of the extracted density was used.
\end{proof}

\paragraph{Resolution and the raw endpoint.}
The signed estimate \eqref{eq:averaged-common-raw-remainder} is a bound
for one common correction tested on unsmoothed windows, not a bound
for its essential supremum or for its absolute integral. The resolution
condition $b_ih_i\to\infty$ cannot hold for a singleton lattice
coordinate when $b_i=o(\sqrt n)$. It therefore does not establish a
fixed-label raw denominator or the microscopic local limit theorem.
The complete complementary integral, the common correction in the
required pointwise norm, and the long-time raw preparation budget
remain the requirements for that same endpoint. Positivity has here
supplied actual probability and signed-average estimates, not an
unproved inference from cutoff transport.
